\section{Conclusion}
\label{sec:conclusion}

We set out to determine whether a large historical corpus can be turned into a
provenance-aware knowledge graph with heavy \mbox{LLM}-agent involvement without
model output becoming indistinguishable from what sources actually say. The answer
this project supports is yes, but only because the constraint was implemented as
mechanism rather than as policy: a model may write exactly one lifecycle status, the
promotion gate is a frozen set of trust classes that does not contain the model's,
the predicate vocabulary is closed with its refusals named, and a cited token absent
from the packet the model was shown is a decidable fabrication.

The resulting artefact is a graph of 165{,}737 nodes and 510{,}905 relationships over
20{,}210 canonical mantras in four named recensions, in which every edge states the
epistemic layer that produced it, the tier that follows from it, the reach of the
claim over its passage, and the textual surface the evidence was read off. Because
the typing is stored rather than documented, the division of labour is measurable:
21.39\,\% of relationships are what a source states, 77.22\,\% are what a
reproducible rule derives from that, 0.35\,\% are interpretive, and \textbf{0.17\,\%
are model-extracted}.

That ratio is the paper's main empirical result, and we want it read correctly. It
does not say agents were unimportant --- they were used continuously across
seventeen days and 249 commits. It says that when the promotion path is closed, the
agentic contribution flows into discovery, measurement, reconciliation and
refutation instead of into assertion. The clearest instance is a gate built to
attack the project's own withholdings rather than its claims, which falsified the
stated reasons of 39 withheld verses using the repository's own released data, and
changed no verse's release status --- only the recorded reason. A method that can
find an error in its own explanations is doing something a method that only validates
its outputs cannot.

The project's failures support the same conclusion more strongly than its successes.
It deleted 41{,}796 of its own relationships in one revision and retired 30{,}539 in
the next; withdrew 96\,\% of its semantic role fillers rather than ship a measured
24.8\,\% error; refused an entire transcription corpus at an inter-reader agreement
of 0.505; refused a semantic-resemblance layer after a re-implemented lexical control
outperformed every candidate; and refused a release candidate while every gate was
green, because completeness turned out to be a property of a registry that had no
terminal state. Twice a refusal's own decisive argument was subsequently found wrong
and withdrawn in turn. Cheap generation makes cheap refusal possible, and refusal is
where the epistemic work happens.

What none of this establishes is the accuracy of the semantic layer. Of 35{,}131
semantic assertions, zero have been reviewed by a human; the strongest review state
in the graph is model adjudication; and on the only independent test available,
single-model adjudication agreed with published human annotation on 75\,\% of
decidable cases. The project says so in its own certification and on its public
pages, and we have been careful to say it here in the same terms. An
evidence-typed graph does not make model output true. It makes model output
\emph{visible}, and it lets a reader who does not trust it exclude it with a
predicate --- which, for scholarship over a corpus like this one, is the more useful
property.
